\section{From $\opto_\varepsilon(K)$ oracles to $\ListLearn$}
\label{sec:listlearn}



We first define the notion of well-separatedness. 
\begin{defn}
    We say that a polytope $K$ with vertex set $V$ is {\em $\delta$-well-separated} if for every vertex $v \in V$, we have
    $$ \dst(v,\CH(V \setminus \{v\})) \geq \delta \cdot \Delta(K),$$
    where $\Delta(K)$ denotes the diameter of $K$.
\end{defn}

We first show the lower bound result~\Cref{lower-bound-0}, which is restated here. 

\lowerbound*

\begin{proof}
    The proof is by producing an ``adversarial oracle''. The parameter $\varepsilon$ is set to $8 \ln d/\sqrt{d}$. Our $K$ will a ``needle'' of the form $\{\lambda u:\lambda\in [-1,1]\}$, where, $u$ is a unit length vector. The oracle's answer for each query will  always be the 0 vector. Let  $v_1,v_2,\ldots, v_q,$ where $q=d^c$ for a constant $c$, be the vectors on which the oracle is queried. The answer 0 is clearly a valid answer for any $\opto_\varepsilon(K)$ oracle,  provided  $u$ is a unit vector with $|u\cdot v_i|\leq 4\ln d/\sqrt d$. For a single $v_i$, the probability that a random $u$ satisfies $|u\cdot v_i|>4\ln d/\sqrt d$ is at most $e^{-d\varepsilon^2}$. So by union bound, we see that there is a vector $u_1$ satisfying 
    $$|u_1\cdot v_i|\leq 4\ln d/\sqrt d\forall i\in [q].$$
    Further, by a similar argument, there is a vector $u_2$satisfying
    $$|u_2\cdot v_i|\leq 4\ln d/\sqrt d, i\in [q]\; ;\; |u_1-u_2|,|u_1+u_2|\geq .1.$$ It follows that the two possible $K$, where 
    $$\{ \lambda u_1:\lambda\in [-1,1]\}\; ,\; \{ \lambda u_2:\lambda\in [-1,1]\}$$
are both consistent with the answer 0 for all $v_i,i\in [q]$ and in addition, we have that no point $w\in {\bf R}^d$ is within distance $\varepsilon$ of two vertices, one from each of these 2 needles. So no answer is valid for both needles and the adversarial oracle can choose one of the needles to render the algorithm's answer incorrect.
\end{proof}

We now formally prove~\Cref{list-0}, which gives an algorithm for the $\ListLearn$ problem. 


\begin{theorem}\label{upper-bound-1}
Suppose $\varepsilon,\delta$ are reals in $[0,1]$ with $\delta^2\geq c\varepsilon\sqrt d, \delta^3 \geq c \, \varepsilon,$ where $c$ is a large enough constant.  Let $K$ be a $\delta$-well-separated $k$-vertex polytope. Suppose we are also given  $ \opto_\varepsilon(K)$ oracle $\cal O$. 
Let $W$ be the set of answers of the oracle to $m=poly(d)\cdot k^{\Omega(1/\delta^2)}$ independent random queries. Then for each vertex $v$ of $K$, there is a point $v' \in W$ such that $|v-v'| \leq O(\delta^2 \Delta(K)/c).$
\end{theorem}
\begin{proof}
 The algorithm chooses a set $U$ of $k^{\Omega(1/\delta^2)}$ unit length i.i.d. Gaussian vectors. For each $u\in U$, it calls the oracle $\cal O$ to find a vector $x(u)$. Let $W$ denote the set $\{x(u): u \in U\}$. We invoke~\Cref{cor:soft} with the parameters: 
$$ \delta' := \delta/4, \, \varepsilon' := \frac{32 \delta^2}{c}$$
on the set $W$. The algorithm in that result outputs a set $Q$. We output that $Q$ as the approximation to the set of vertices of $K$. We now prove that the set $Q$ has the desired properties. We first show that for every vertex of $K$, there is a direction $u$ in $U$ along which the projection of this vertex is higher than the projection of the remaining vertices by a large enough margin. 
Let $\Mdot{1}, \ldots, \Mdot{k}$ be the vertices of $K$. 
\begin{claim}
    \label{cl:ul}
    With high probability, the following event happens: for each $\ell \in [k]$, there is a vector $u^{(\ell)} \in U$ such that for all $\ell' \in [k], \ell' \neq \ell,$ we have
    \begin{align}
        \label{eq:100}
        u^{(\ell)}\cdot M_{\cdot, \ell}> u^{(\ell)}\cdot M_{\cdot, \ell'}+\frac{c \, \varepsilon \cdot \Delta(K)}{8 \delta^2}
    \end{align}
\end{claim}

\begin{proof}
    Fix a vertex $\Mdot{\ell}$ of $K$.  Let $K'$ be the convex hull of $\{\Mdot{1}, \ldots, \Mdot{k} \} \setminus \{\Mdot{\ell}\}.$ We invoke~\Cref{thm:sephp} on the polytope $K'$ and the point $a := \Mdot{\ell}$. The definition of $\delta$-well-separated implies that~\eqref{eq:condhp} is satisfied with $\rho = \delta$. Using $V = \Re^d$ in the statement of~\Cref{thm:sephp}, we see that $$\Pr_{u \in U} \left[ u \cdot \Mdot{\ell} - \max_{\ell' \in [k], \ell' \neq \ell} u \cdot \Mdot{\ell'} \geq \frac{\delta \sqrt{\log k} \Delta(K)}{\sqrt{\log k} + 4 \delta \sqrt{d}} \right] \geq \frac{1}{40} k^{-10/\delta^2}.$$
    Since 
    $$\frac{\delta \sqrt{\log k}}{\sqrt{\log k} + 4 \delta \sqrt{d}}  \geq \min \left( \frac{\delta}{2}, \frac{1}{8 \sqrt{d}} \right) \geq \frac{\varepsilon c}{8 \delta^2}, $$
    the desired result follows from the fact that $|U| \gg \frac{1}{40} k^{-10/\delta^2}.$
    
\end{proof}

For rest of the proof, assume that the statement in~\Cref{cl:ul} holds true, i.e., there are directions $u^{(1)}, \ldots, u^{(k)} \in U$ satisfying~\eqref{eq:100}. 
We now show that for every vertex of $\Mdot{\ell}$ of $K$, the corresponding point $x(u^{\ell})$ is close to  $\Mdot{\ell}$. 
\begin{claim}
    \label{cl:xl}
    For every $\ell \in [k]$, 
    $$|x(u^{(\ell)})-M_{\cdot,\ell}|\leq 17 \delta^2 \Delta(K)/c.$$
\end{claim}
\begin{proof}
    By the definition of $\cal O$, we know that $x(u^{(\ell)})$ can be written as $y(u^{\ell})+z(\uell)$, where $y(\uell) \in K$ and $|z(\uell)| \leq \varepsilon \Delta(K)$. Thus, there is a convex combination $\lambda_{\ell'}, \ell' \in [k],$ of the vertices $\Mdot{\ell'}$ of $K$ such that 
    $$ x(\uell) = \sum_{\ell' \in [k]} \lambda_{\ell'} \Mdot{\ell'} + z(\uell).$$
   

By the definition of $\cal O$, $x(u^{(\ell)})\cdot u^{(\ell)}\geq M_{\cdot, \ell}\cdot u^{(\ell)}-\varepsilon\Delta (K)$ and $|z(u^{(\ell)})\cdot u^{(\ell)}|\leq\varepsilon\Delta(K)$. So, we get
$$M_{\cdot,\ell}\cdot u^{(\ell)}-\varepsilon\Delta(K)\leq \sum_{\ell' \in [k]}\lambda_{\ell'} M_{\cdot,\ell'}\cdot u^{(\ell')}+\varepsilon\Delta(K),$$
which implies (after subtracting $\lambda_\ell M_{\cdot,\ell} \cdot u^{(\ell)}$ from both sides):
$$(1-\lambda_\ell) M_{\cdot,\ell}u^{(\ell)}\leq \sum_{\ell'\not= \ell} \lambda_{\ell'} M_{\cdot,\ell'}\cdot u^{(\ell')}+2\varepsilon\Delta(K)$$
which, using~\Cref{cl:ul}, yields:
$$(1-\lambda_\ell) M_{\cdot,\ell}u^{(\ell)}\leq (1-\lambda_\ell) M_{\cdot,\ell}u^{(\ell)}-(1-\lambda_\ell)\frac{c\varepsilon\Delta(K)}{8\delta^2}+2\varepsilon\Delta(K).$$
%which leads to:


    
It follows from the above inequality that 
\begin{align}
\label{eq:105}
    1- \lambda_\ell \leq \frac{16 \delta^2}{c}.
\end{align}
    Therefore, 
\begin{align}
    |x(\uell) - \Mdot{\ell}| = \left| \sum_{\ell' \neq \ell} \lambda_{\ell'} (\Mdot{\ell} - \Mdot{\ell'})  
    \right| + \varepsilon  \Delta(K) \leq \sum_{\ell' \neq \ell} \lambda_{\ell'} \Delta(K) + \varepsilon \Delta(K) \stackrel{\eqref{eq:105}}{\leq} \frac{17 \delta^2 \Delta(K)}{c}.
    \notag
\end{align}
%where the second inequality follows from the fact that $\Delta(K) \leq 1$. 
\end{proof}

This completes proof of the Theorem. 
\end{proof}

The statement of~\Cref{upper-bound-1} requires two different bounds relating $\delta$ to $\varepsilon$. In some applications, it may be difficult to ensure both of these conditions. We consider the setting when $d \gg k$, where the following variation of~\Cref{upper-bound-1} is better suited. 

\begin{theorem}\label{upper-bound-2}
Suppose $\varepsilon,\delta$ are reals in $[0,1]$ with $\delta^2\geq c\varepsilon\sqrt d, \delta \geq \frac{ \sqrt{\log k}}{\sqrt{cd}},$ where $c$ is  large enough constant.  Let $K$ be a $\delta$-well-separated $k$-vertex polytope. Suppose we are also given an  $\opto_\varepsilon(K)$ oracle $\cal O$.
Let $W$ be the set of answers of the oracle to $m=poly(d)\cdot k^{O(1/\delta^2)}$ independent queries.  Then for each vertex $v$ of $K$, there is a point $v' \in W$ such that $|v-v'| \leq O(\delta^2 \Delta(K)/\sqrt{c}).$
\end{theorem}
\begin{proof}
    The proof is identical to that of~\Cref{upper-bound-1}, except that in the proof of~\Cref{cl:ul}, we now have the following modified argument. Since $\delta \sqrt{cd} \geq \sqrt{\log k}, $ we get 
     $$\frac{\delta \sqrt{\log k}}{\sqrt{\log k} + 4 \delta \sqrt{d}}  \geq \frac{1}{8 \sqrt{cd}} \geq \frac{\varepsilon \sqrt{c}}{8 \delta^2}. $$
     Rest of the proof of~\Cref{cl:ul} follows without any further changes (where we replace $c$ by $\sqrt{c}$). 
\end{proof}
